% Options for packages loaded elsewhere
\PassOptionsToPackage{unicode}{hyperref}
\PassOptionsToPackage{hyphens}{url}
\documentclass[
  a4paper,11pt]{article}
\usepackage{xcolor}
\usepackage[margin=1in]{geometry}
\usepackage{amsmath,amssymb}
\usepackage{graphicx}
\setcounter{secnumdepth}{-\maxdimen} % remove section numbering
\usepackage{iftex}
\ifPDFTeX
  \usepackage[T1]{fontenc}
  \usepackage[utf8]{inputenc}
  \usepackage{textcomp}
\else
  \usepackage{unicode-math}
  \defaultfontfeatures{Scale=MatchLowercase}
  \defaultfontfeatures[\rmfamily]{Ligatures=TeX,Scale=1}
\fi
\usepackage{lmodern}
\ifPDFTeX\else
  \setmainfont[]{Times New Roman}
\fi
\IfFileExists{upquote.sty}{\usepackage{upquote}}{}
\IfFileExists{microtype.sty}{%
  \usepackage[]{microtype}
  \UseMicrotypeSet[protrusion]{basicmath}
}{}
\makeatletter
\@ifundefined{KOMAClassName}{%
  \IfFileExists{parskip.sty}{%
    \usepackage{parskip}
  }{%
    \setlength{\parindent}{0pt}
    \setlength{\parskip}{6pt plus 2pt minus 1pt}}
}{% if KOMA class
  \KOMAoptions{parskip=half}}
\makeatother
\usepackage{longtable,booktabs,array}
\newcounter{none}
\usepackage{calc}
\usepackage{etoolbox}
\makeatletter
\patchcmd\longtable{\par}{\if@noskipsec\mbox{}\fi\par}{}{}
\makeatother
\IfFileExists{footnotehyper.sty}{\usepackage{footnotehyper}}{\usepackage{footnote}}
\makesavenoteenv{longtable}
\usepackage{bookmark}
\IfFileExists{xurl.sty}{\usepackage{xurl}}{}
\urlstyle{same}
\usepackage{hyperref}
\hypersetup{
  hidelinks,
  pdfcreator={LaTeX via pandoc}}
\setlength{\emergencystretch}{3em}
\providecommand{\tightlist}{%
  \setlength{\itemsep}{0pt}\setlength{\parskip}{0pt}}
\usepackage{amsthm}
\newtheorem{theorem}{Theorem}[section]
\newtheorem{proposition}[theorem]{Proposition}
\newtheorem{lemma}[theorem]{Lemma}
\newtheorem{corollary}[theorem]{Corollary}
\theoremstyle{definition}
\newtheorem{definition}[theorem]{Definition}
\theoremstyle{remark}
\newtheorem{remark}[theorem]{Remark}
\newtheorem{observation}[theorem]{Observation}
\begin{document}
\section{Mass Structure of Gauge Bosons --- Higgs Non-Involvement
Indicated by Sign
Vectors}\label{mass-structure-of-gauge-bosons-higgs-non-involvement-indicated-by-sign-vectors}

\subsection{--- A Structural Answer to the W±/Z⁰ Mass Problem Based on
Table A of {[}3{]} and §8 of {[}4{]}
---}\label{a-structural-answer-to-the-wzux2070-mass-problem-based-on-table-a-of-3-and-8-of-4}

\textbf{Author}: Noriaki Kihara (WF System Co., Ltd.)

ORCID: 0009-0002-7389-6873

DOI: https://doi.org/10.5281/zenodo.19731608

\textbf{Date}: April 2026

\begin{center}\rule{0.5\linewidth}{0.5pt}\end{center}

\subsection{§0 Purpose of This Paper}\label{purpose-of-this-paper}

In the Standard Model, the mass acquisition of W± (80.4 GeV) and Z⁰
(91.2 GeV) is explained by the Higgs mechanism (electroweak symmetry
breaking). Within the present framework as well, {[}5{]} §1 lists ``the
mass acquisition mechanism of W±/Z⁰ (details of indirect coupling not
mediated through spatial axes)'' as an open problem.

This paper compares the 6-dimensional sign vectors for each particle
listed in Table A of {[}3{]} with the geometric origin of acceleration
derived in {[}4{]} §8, thereby demonstrating the following:

\begin{enumerate}
\def\labelenumi{\arabic{enumi}.}
\tightlist
\item
  W±/Z⁰ do not share any spatial axis with the Higgs boson and have no
  direct coupling to the Higgs (§1)
\item
  The sign vectors of W±/Z⁰ have the same structure as the photon (zero
  spatial components), and the origin of their ``mass'' is attributable
  not to the Higgs but to the curvature of the subjective space (§2)
\item
  The mechanism of mass acquisition is distinguished into two types: the
  rest mass of fermions (Higgs-derived) and the energy of gauge bosons
  (curvature-derived) (§3)
\end{enumerate}

In this paper, ``Higgs non-involvement'' refers to the absence of the
direct coupling through shared spatial axes as formulated in {[}5{]} §8.
Preliminary considerations on fermion mass hierarchy and quantitative
mass ratios are treated separately in ``Considerations on Mass Structure
--- A Framework for Mass Analysis via Axis Scale Values and Signed
Areas.'' This supplementary lecture addresses the distinct problem
domain of gauge boson mass structure.

\textbf{References}: - {[}1{]} Shape-Invariant Standing Waves on a
Closed Sphere --- Stability by Dimension and the Geometric Necessity of
3+1 Spacetime - {[}2{]} Full-Sphere Coverage of Central Projection in
Discrete Space and the Number-Theoretic Necessity of a 5-Dimensional
Background Space (W5), DOI: 10.5281/zenodo.19624957 - {[}3{]} Set
Structure of the 6-Dimensional Hyperrectangle and Its Combinatorial
Properties (W8), DOI: 10.5281/zenodo.19721125 - {[}4{]} Four Modes of
Shape-Invariant Waves (W10), DOI: 10.5281/zenodo.19709798 - {[}5{]}
Interactions of Shape-Invariant Waves (W11), DOI:
10.5281/zenodo.19721128

\begin{center}\rule{0.5\linewidth}{0.5pt}\end{center}

\subsection{§1 Comparison of Sign
Vectors}\label{comparison-of-sign-vectors}

\subsubsection{1.1 Higgs Boson and Gauge
Bosons}\label{higgs-boson-and-gauge-bosons}

Based on Table A of {[}3{]}, we enumerate the 6-dimensional sign vectors
\((x, y, z, t, R, Q)\) of the Higgs boson and spin-1 gauge bosons.

{\def\LTcaptype{none} % do not increment counter
\begin{longtable}[]{@{}
  >{\raggedright\arraybackslash}p{(\linewidth - 8\tabcolsep) * \real{0.1905}}
  >{\centering\arraybackslash}p{(\linewidth - 8\tabcolsep) * \real{0.2381}}
  >{\raggedright\arraybackslash}p{(\linewidth - 8\tabcolsep) * \real{0.1905}}
  >{\raggedright\arraybackslash}p{(\linewidth - 8\tabcolsep) * \real{0.1905}}
  >{\raggedleft\arraybackslash}p{(\linewidth - 8\tabcolsep) * \real{0.1905}}@{}}
\toprule\noalign{}
\begin{minipage}[b]{\linewidth}\raggedright
Particle
\end{minipage} & \begin{minipage}[b]{\linewidth}\centering
Spin
\end{minipage} & \begin{minipage}[b]{\linewidth}\raggedright
Sign vector \((x, y, z, t, R, Q)\)
\end{minipage} & \begin{minipage}[b]{\linewidth}\raggedright
Non-zero axes
\end{minipage} & \begin{minipage}[b]{\linewidth}\raggedleft
Measured mass
\end{minipage} \\
\midrule\noalign{}
\endhead
\bottomrule\noalign{}
\endlastfoot
Higgs \(H^0\) & 0 & \((0, 0, +, 0, 0, 0)\) & z & 125.3 GeV \\
Photon \(\gamma\) & 1 & \((0, 0, 0, 0, +, 0)\) & R & 0 \\
\(Z^0\) & 1 & \((0, 0, 0, 0, -, 0)\) & R & 91.2 GeV \\
\(W^+\) & 1 & \((0, 0, 0, +, 0, 0)\) & t & 80.4 GeV \\
\(W^-\) & 1 & \((0, 0, 0, -, 0, 0)\) & t & 80.4 GeV \\
\end{longtable}
}

\subsubsection{1.2 Non-Sharing of Spatial
Axes}\label{non-sharing-of-spatial-axes}

\textbf{Observation 1.1}. The only non-zero axis of the Higgs boson
\(H^0 = (0, 0, +, 0, 0, 0)\) is the spatial \(z\) axis. In contrast, the
photon, \(Z^0\), and \(W^{\pm}\) all have \textbf{spatial components
\((x, y, z)\) entirely equal to zero}.

Therefore, the photon, \(Z^0\), and \(W^{\pm}\) do not share any spatial
axis with the Higgs. The mass acquisition mechanism formulated in
{[}5{]} §8 (phase coupling on shared spatial axes with standing waves)
does not apply to these gauge bosons.

\textbf{Proposition 1.1}. The masses of \(W^{\pm}\) and \(Z^0\) cannot
be explained by rest mass acquisition through shared spatial axes with
the Higgs boson ({[}5{]} §8). An explanation via a different mechanism
is required; in §2, it is described as acceleration energy originating
from the curvature of central projection.

\emph{Argument}. By {[}5{]} §8.1, coupling with a standing wave (Higgs,
\(\kappa = 1\)) occurs when a direction-constrained wave
(\(\kappa = 2\)) directly shares a spatial axis with the standing wave.
\(W^{\pm}\) (non-zero only on the \(t\) axis) and \(Z^0\) (non-zero only
on the \(R\) axis) have all spatial components equal to zero, and thus
share no spatial axis with the non-zero axis (\(z\)) of the Higgs.
\(\square\)

\subsubsection{1.3 Structural Identity with the
Photon}\label{structural-identity-with-the-photon}

\textbf{Observation 1.2}. \(W^{\pm}\), \(Z^0\), and the photon
\(\gamma\) all have zero spatial components. Their structure with
respect to coupling with the Higgs is identical among the three.

{\def\LTcaptype{none} % do not increment counter
\begin{longtable}[]{@{}lccc@{}}
\toprule\noalign{}
Property & \(\gamma\) & \(Z^0\) & \(W^{\pm}\) \\
\midrule\noalign{}
\endhead
\bottomrule\noalign{}
\endlastfoot
Spatial components \((x,y,z)\) & \((0,0,0)\) & \((0,0,0)\) &
\((0,0,0)\) \\
Shared spatial axis with Higgs & None & None & None \\
Rest mass acquisition via Higgs & None & None & None \\
\end{longtable}
}

It is well known that the photon ``has zero rest mass yet possesses
energy.'' \(W^{\pm}\) and \(Z^0\) share the same Higgs-non-involved
structure as the photon.

\begin{center}\rule{0.5\linewidth}{0.5pt}\end{center}

\subsection{§2 Origin of Gauge Boson
Energy}\label{origin-of-gauge-boson-energy}

\subsubsection{2.1 Sign Vector of the
Graviton}\label{sign-vector-of-the-graviton}

In Table A of {[}3{]}, the sign vector of the graviton (spin 2, \(tR\)
type) is:

\[G = (0, 0, 0, \pm, \pm, 0)\]

The graviton has non-zero components on both the \(t\) axis and the
\(R\) axis. \(W^{\pm}\) has a non-zero component only on the \(t\) axis,
while \(Z^0\) and \(\gamma\) have non-zero components only on the \(R\)
axis; in each case, their non-zero axes are subsets of those of the
graviton. That is, the non-zero axes \(t, R\) of \(W^{\pm}\) and \(Z^0\)
share the same axis structure as the acceleration mechanism mediated by
the graviton ({[}4{]} §8, detailed in §2.2). Therefore, the energy of
these bosons can also be described as curvature-derived acceleration in
accordance with {[}4{]} §8.

\subsubsection{2.2 Geometric Origin of
Acceleration}\label{geometric-origin-of-acceleration}

By {[}4{]} §8, the acceleration of shape-invariant waves is
intrinsically derived from the geometry of central projection:

\begin{enumerate}
\def\labelenumi{(\roman{enumi})}
\item
  \textbf{Global acceleration} ({[}4{]} §8.2): The curvature radius
  \(R\) of \(S^3\) imparts a global acceleration
  \(\Delta^2\mathbf{k}_{\text{global}} = f(R, \mathbf{k}_t)\) to wave
  trajectories
\item
  \textbf{Local acceleration} ({[}4{]} §8.3): The Hessian \(D^2\pi\) of
  the subspace central projection determines a local acceleration
  \(\Delta^2\mathbf{k}_{\text{local}}\) at each point
\item
  \textbf{Unification} ({[}4{]} §8.4 Proposition 8.2): Global curvature
  and local nonlinearity are manifestations at different scales of a
  single map, the \(R\)-parameterized central projection
  \(\pi: S^5(R) \to \mathbb{Z}^4\)
\end{enumerate}

What is decisive here is that the acceleration function
\(f(R, \mathbf{k}_t)\) depends on the value of \(R\) (the curvature
radius). The central projection map itself is defined as the same
regular map regardless of the sign of \(R\) (as proved in {[}2{]}), but
the value of the acceleration function differs according to the sign of
\(R\), and this manifests as a difference in acceleration energy.

\subsubsection{2.3 Relationship Between R-Axis Sign and
Energy}\label{relationship-between-r-axis-sign-and-energy}

\textbf{Proposition 2.1}. The energy difference between the photon
(\(R = +\)) and \(Z^0\) (\(R = -\)) originates from the difference in
sign of the curvature radius \(R\) of the subjective space.

\emph{Argument}. By {[}4{]} §8.2, acceleration is given by the function
\(f(R, \mathbf{k}_t)\) of the curvature radius \(R\). The central
projection \(\pi: S^5(R) \to \mathbb{Z}^4\) is defined as the same
regular map regardless of the sign of \(R\) (as proved in {[}2{]}), but
the sign of \(R\) itself governs the value of the acceleration function
\(f(R, \mathbf{k}_t)\). For \(R = +\) (photon) and \(R = -\) (\(Z^0\)),
different acceleration energies are measured, and this manifests as the
energy difference between the two.

From the perspective of an external observer (a different subjective
space), the energy of a wave localized in a subjective space with
\(R < 0\) includes acceleration energy originating from the curvature of
that space. This is measured as the 91.2 GeV of \(Z^0\). \(\square\)

\subsubsection{2.4 t-Axis Sign and the Energy of
W±}\label{t-axis-sign-and-the-energy-of-w}

The same logic applies to \(W^{\pm}\) (non-zero only on the \(t\) axis).
The structure along the \(t\)-axis direction determines a specific
energy scale, which is measured as 80.4 GeV.

\begin{center}\rule{0.5\linewidth}{0.5pt}\end{center}

\subsection{§3 Two Types of Mass
Acquisition}\label{two-types-of-mass-acquisition}

\subsubsection{3.1 Distinction of Types}\label{distinction-of-types}

Within the present framework, what is measured as ``mass'' has two
structurally distinct origins.

\textbf{Type I: Higgs-Derived Rest Mass} (Fermions)

\begin{itemize}
\tightlist
\item
  Mechanism: Phase coupling on shared spatial axes with standing waves
  (\(\kappa = 1\), Higgs) ({[}5{]} §8)
\item
  Target: Direction-constrained waves (\(\kappa = 2\)) with non-zero
  spatial axes, namely fermions
\item
  Result: Velocity limitation due to reduction of effective spatial
  components \(k_{s_0}^{\text{eff}} = k_{s_0} - \delta k_H\) ({[}5{]}
  §8.2)
\item
  Characteristic: The generational mass hierarchy is determined by the
  presence or absence of shared spatial axes with the Higgs
\end{itemize}

\textbf{Type II: Curvature-Derived Energy} (Gauge bosons with zero
spatial components)

Type II gauge bosons have zero spatial components and therefore do not
possess rest mass through shared spatial axes with the Higgs. The
observed ``mass'' is described within the present framework as
acceleration energy originating from the curvature of the subjective
space.

\begin{itemize}
\tightlist
\item
  Mechanism: Acceleration dependent on the curvature radius \(R\) of the
  subjective space ({[}4{]} §8). The central projection is defined as
  the same regular map regardless of the sign of \(R\) ({[}2{]}), but
  the value of the acceleration function depends on the sign of \(R\)
\item
  Target: Bosons with zero spatial components (\(\gamma\), \(Z^0\),
  \(W^{\pm}\))
\item
  Result: Acceleration energy determined by the acceleration function
  \(f(R, \mathbf{k}_t)\)
\item
  Characteristic: The signs of the \(R\) axis and \(t\) axis are
  reflected in the energy
\end{itemize}

\subsubsection{3.2 Classification Table}\label{classification-table}

{\def\LTcaptype{none} % do not increment counter
\begin{longtable}[]{@{}
  >{\raggedright\arraybackslash}p{(\linewidth - 4\tabcolsep) * \real{0.3333}}
  >{\raggedright\arraybackslash}p{(\linewidth - 4\tabcolsep) * \real{0.3333}}
  >{\raggedright\arraybackslash}p{(\linewidth - 4\tabcolsep) * \real{0.3333}}@{}}
\toprule\noalign{}
\begin{minipage}[b]{\linewidth}\raggedright
Property
\end{minipage} & \begin{minipage}[b]{\linewidth}\raggedright
Type I (Higgs-derived)
\end{minipage} & \begin{minipage}[b]{\linewidth}\raggedright
Type II (Curvature-derived)
\end{minipage} \\
\midrule\noalign{}
\endhead
\bottomrule\noalign{}
\endlastfoot
Target & Fermions (\(\kappa = 2\)) & Gauge bosons
(\(\gamma, Z^0, W^{\pm}\)) \\
Spatial components & Non-zero (1 or 2 axes) & Zero \\
Coupling with Higgs & Direct coupling via shared spatial axes & None \\
Mechanism & Velocity limitation \(v^{\text{eff}} < c\) ({[}5{]} §8.2) &
Curvature acceleration ({[}4{]} §8) \\
Generation dependence & Present (\(m_3 \gg m_{1,2}\)) & Absent \\
Structural difference from photon & Present (presence/absence of spatial
components) & Absent (identical structure) \\
\end{longtable}
}

\subsubsection{3.3 Comparison with the Standard
Model}\label{comparison-with-the-standard-model}

In the Standard Model, both fermion masses (Yukawa coupling) and gauge
boson masses (covariant derivative of the Higgs field) are explained in
a unified manner by the Higgs mechanism. In the present framework, the
two types separate naturally from the structure of sign vectors.

{\def\LTcaptype{none} % do not increment counter
\begin{longtable}[]{@{}
  >{\raggedright\arraybackslash}p{(\linewidth - 4\tabcolsep) * \real{0.3333}}
  >{\raggedright\arraybackslash}p{(\linewidth - 4\tabcolsep) * \real{0.3333}}
  >{\raggedright\arraybackslash}p{(\linewidth - 4\tabcolsep) * \real{0.3333}}@{}}
\toprule\noalign{}
\begin{minipage}[b]{\linewidth}\raggedright
Particle group
\end{minipage} & \begin{minipage}[b]{\linewidth}\raggedright
Standard Model mechanism
\end{minipage} & \begin{minipage}[b]{\linewidth}\raggedright
Type in this framework
\end{minipage} \\
\midrule\noalign{}
\endhead
\bottomrule\noalign{}
\endlastfoot
Fermions & Yukawa coupling \(y_f \langle\phi\rangle\) & Type I (shared
spatial axes) \\
\(W^{\pm}, Z^0\) & Covariant derivative \(|D_\mu\phi|^2\) & Type II
(curvature energy) \\
Photon \(\gamma\) & No coupling (zero mass) & Type II (curvature energy,
\(R > 0\)) \\
Gluon & No coupling (zero mass) & --- (\(Q\)-axis transition type,
separate mechanism) \\
\end{longtable}
}

Gluons have zero spatial components, but their non-zero axis involves
discrete label transitions on the \(Q\) axis ({[}5{]} §5), which is
qualitatively different from the \(t, R\) axis structure of the photon
and \(W^{\pm}/Z^0\). Therefore, they do not belong to Type II either and
are treated as a separate mechanism.

In the Standard Model, the mass difference between \(W^{\pm}/Z^0\) and
the photon is explained by the vacuum expectation value of the Higgs
field and electroweak symmetry breaking, whereas in the present
framework, this difference reduces to the difference in acceleration
energy determined by the signs of the \(R\) axis and \(t\) axis.

\textbf{Remark}: In the Standard Model, \(W^{\pm}/Z^0\) are treated as
massive gauge bosons possessing rest mass, but in the present framework
they have no rest mass because their spatial components are zero; the
measured 80.4 GeV and 91.2 GeV are measured as curvature-derived energy.
This difference in interpretation reflects the fact that the definition
of ``mass'' in the present framework differs from that in the Standard
Model. A detailed examination of the relationship with the Standard
Model's electroweak precision measurements (\(W/Z\) mass ratio,
\(\sin^2\theta_W\), etc.) lies outside the scope of this paper.

\begin{center}\rule{0.5\linewidth}{0.5pt}\end{center}

\subsection{§4 Conclusions}\label{conclusions}

In this paper, we have provided the following structural answers to the
mass acquisition mechanism of W±/Z⁰, which was listed as an open problem
in {[}5{]} §1.

\begin{enumerate}
\def\labelenumi{\arabic{enumi}.}
\item
  \textbf{W±/Z⁰ do not share any spatial axis with the Higgs}
  (Proposition 1.1). This can be confirmed directly from the sign
  vectors in Table A of {[}3{]}, and the mass acquisition mechanism of
  {[}5{]} §8 (phase coupling via shared spatial axes) does not apply.
\item
  \textbf{W±/Z⁰ have the same Higgs-non-involved structure as the
  photon} (Observation 1.2). The structural feature that spatial
  components are zero is common to all three.
\item
  \textbf{Mass acquisition is distinguished into two types}. The rest
  mass of fermions (Type I: shared spatial axes with the Higgs) and the
  energy of gauge bosons (Type II: acceleration energy originating from
  the curvature of the subjective space, {[}4{]} §8).
\item
  \textbf{The energy difference between the photon and Z⁰ reduces to the
  difference in R-axis sign} (Proposition 2.1). The central projection
  is defined as the same regular map regardless of the sign of \(R\)
  ({[}2{]}), but the sign of \(R\) governs the value of the acceleration
  function; different energies are measured for \(R = +\) (photon) and
  \(R = -\) (\(Z^0\)).
\item
  \textbf{The two types of mass acquisition separate naturally within
  the present framework} (§3). The distinction between fermions (Type I)
  and gauge bosons (Type II) follows immediately from the presence or
  absence of spatial components in the sign vectors. In the Standard
  Model, the Higgs mechanism explains both in a unified manner, but in
  the present framework, each has a different geometric origin.
\end{enumerate}

The derivation of specific numerical values for the energy of the
photon, \(Z^0\), and \(W^{\pm}\) (the quantitative relationship between
the sign and curvature radius of \(R\)) lies outside the scope of this
paper and remains an open problem.

\begin{center}\rule{0.5\linewidth}{0.5pt}\end{center}

\subsection{References}\label{references}

\begin{itemize}
\tightlist
\item
  {[}1{]} Noriaki Kihara, ``Shape-Invariant Standing Waves on a Closed
  Sphere --- Stability by Dimension and the Geometric Necessity of 3+1
  Spacetime,'' 2026.
\item
  {[}2{]} Noriaki Kihara, ``Full-Sphere Coverage of Central Projection
  in Discrete Space and the Number-Theoretic Necessity of a
  5-Dimensional Background Space,'' Zenodo preprint, DOI:
  10.5281/zenodo.19624957 (2026).
\item
  {[}3{]} Noriaki Kihara, ``Set Structure of the 6-Dimensional
  Hyperrectangle and Its Combinatorial Properties,'' Zenodo preprint,
  DOI: 10.5281/zenodo.19721125 (2026).
\item
  {[}4{]} Noriaki Kihara, ``Four Modes of Shape-Invariant Waves ---
  Spin, Chirality, and Statistics Determined by Wave-Vector Structure,''
  Zenodo preprint, DOI: 10.5281/zenodo.19709798 (2026).
\item
  {[}5{]} Noriaki Kihara, ``Interactions of Shape-Invariant Waves ---
  Axial Displacement Transfer, Wave-Packet Deformation, and Retroactive
  Constitution of Causality,'' Zenodo preprint, DOI:
  10.5281/zenodo.19721128 (2026).
\end{itemize}
\end{document}
